\newcommand{\ponisti}[2]{\tikz[baseline=(x.base)]{\node[inner sep=0pt] (x) {$#2$};\draw[#1,line width=.3pt] (x.south west)--(x.north east);}}
\slajdid{04-s004}
\begin{frame}[label=04-s004]
\gusttekst
Bez obzira na osnovu pozicionog brojnog sistema uvek je $RV<1$
\[RV=\sum_{i=-k}^{-1}c_i r^i\qquad RV_{max}=\sum_{i=-k}^{-1}(c_i)_{max}r^i\]
\[c_i\in\{0,1,2,\ldots r-2,r-1\}\qquad(c_i)_{max}=r-1\]
\[RV_{max}=\sum_{i=-k}^{-1}(r-1)r^i=\sum_{i=-k}^{-1}\left(r^{i+1}-r^i\right)\]
% Iste oznake poništavanja kao u originalu, izrađene kao vektorske linije.
\[=\left(r^0-\ponisti{orange}{r^{-1}}\right)+\left(\ponisti{orange}{r^{-1}}-\ponisti{blue}{r^{-2}}\right)+\cdots+\left(\ponisti{green!60!black}{r^{-k+2}}-\ponisti{black}{r^{-k+1}}\right)+\left(\ponisti{black}{r^{-k+1}}-r^{-k}\right)\]
\[RV_{max}=\left(1-r^{-k}\right)<1\]
Koliko će biti „blizu 1“ zavisi od broja upotrebljenih razlomljenih mesta za prikaz broja\\
(normalno i od osnove ali razmišljamo o istoj osnovi)\\
NAPOMENA: Ovo važi za brojne sisteme sa celobrojnim osnovama
\end{frame}
